\section{Installing and launching}
\label{sec:launching}

\subsection{Installation}

Three ways to install, ordered by convenience. Requires
Python 3.10 or newer.

\subsubsection*{Option A: From PyPI (recommended)}

\begin{lstlisting}[language=bash]
pip install al-dic
\end{lstlisting}

This pulls the latest released wheel and every runtime dependency
in one step.

\subsubsection*{Option B: From a GitHub Release wheel}

Useful when you cannot reach PyPI (corporate firewall, air-gapped
network) or when you want to install a specific past version whose
artefacts are still attached to its release page.

\begin{enumerate}
  \item Go to \url{https://github.com/zachtong/pyALDIC/releases}.
  \item Expand \emph{Assets} on the release you want; download the
    \texttt{al\_dic-<version>-py3-none-any.whl} file.
  \item Install locally:
\begin{lstlisting}[language=bash]
pip install ./al_dic-<version>-py3-none-any.whl
\end{lstlisting}
\end{enumerate}

Dependencies are still fetched from PyPI at install time, so this
method only bypasses PyPI for the \texttt{al-dic} package itself.

\subsubsection*{Option C: From source (development)}

Editable install that lets you modify the code and see the effect
without re-installing.

\begin{lstlisting}[language=bash]
git clone https://github.com/zachtong/pyALDIC.git
cd pyALDIC
pip install -e ".[dev]"
\end{lstlisting}

The \texttt{.[dev]} extras include \texttt{pytest} and
\texttt{pytest-xdist} for running the test suite.

\subsection{Launching the GUI}

\begin{lstlisting}[language=bash]
al-dic
# or equivalently
python -m al_dic
\end{lstlisting}

The window opens at a fixed 1380$\times$800 minimum size. You can
resize it to taste; the three columns keep their proportions.

\subsection{What you should see first}

The first time the window opens:
\begin{itemize}
  \item The left image list is empty and shows a drop zone that reads
    \emph{``Drop image folder or Browse''}.
  \item The centre canvas shows an empty dark background.
  \item The right sidebar shows a disabled \textbf{Run DIC Analysis}
    button (you cannot run without images and a Region of Interest).
\end{itemize}

\begin{tipbox}
All user-visible text uses the full phrase \emph{Region of Interest}
instead of the DIC jargon ``ROI''. The 50-pixel Region column in the
image list is abbreviated to ``Region'' to fit, with
Need/Edit/Add button labels inside.
\end{tipbox}
